%% Supplement S2: re-signing an archive (moved from the paper).
Section~\ref{sec:renewal} of the paper sets out what a renewal must preserve.
This section measures the computation it shares with every reading of that
requirement, hashing the asset and signing a new claim, and leaves out
validating the original evidence, serialising it as an ingredient, obtaining a
new timestamp and writing the result.
\subsection{A model with no fitted parameters}

Per asset the work splits in two. Hashing the bytes for the hard binding
scales with the file; building the claim, signing it and assembling the
manifest does not. So
\renewcommand{\theequation}{S\arabic{equation}}
\begin{equation}
t(s) = a + s/h
\label{eq:resign}
\end{equation}
with $s$ the file size, $h$ the hash throughput and $a$ the fixed cost. Both
terms are measured separately and neither is fitted to the data the model is
tested on. Fitting both by least squares fails instructively: signing time
through OpenSSL varies by more than the whole size-dependent term at a hundred
kilobytes, so the regression absorbs signing noise into the slope and reports
the same hash running eight times slower under ML-DSA-87 than under ECDSA.

\begin{table}[t]
\caption{The per-asset model and its test. $a$ is the fixed cost, measured over a digest already in hand; the hash term is 2,417~MB/s measured with no signing in the loop. Neither is fitted. The error column is the model's median error against held-out images it was not built from. \emph{Crossover} is the file size at which hashing costs as much as everything else, which is where the algorithm stops mattering.}
\label{tab:archivemodel}
\begin{tabular}{lrrrrr}
\toprule
Algorithm & $a$ (ms) & mean/median & $a$ liboqs (ms) & Held-out error & Crossover\\
\midrule
ES256 & 0.031 & 1.01 & -- & 0.24\% & 73~KiB\\
ML-DSA-44 & 0.293 & 1.35 & 0.087 & 10.28\% & 692~KiB\\
ML-DSA-65 & 0.502 & 1.16 & 0.129 & 2.08\% & 1,186~KiB\\
ML-DSA-87 & 0.596 & 1.22 & 0.183 & 0.43\% & 1,407~KiB\\
\bottomrule
\end{tabular}
\end{table}


Table~\ref{tab:archivemodel} reports the model and its held-out test. The
crossover column is the operational result, the file size at which hashing
costs as much as everything else. At ES256 it is \NarchCrossES\,KiB, so almost
every stored photograph is hash-bound today; under ML-DSA it moves to between
\NarchCrossLo{} and \NarchCrossHi\,MiB, past a web image and short of a camera
JPEG. Through the unoptimised build of Section~\ref{sec:supp-impl} it lands past a camera JPEG and would tell an operator the opposite of the truth about a camera archive. The held-out error is largest for ML-DSA-44, \NarchHeldDSAfortyfour\% against \NarchHeldDSAsixtyfive{} and \NarchHeldDSAeightyseven\% for the larger sets, and it follows the spread of the fixed cost rather than the parameter set: in our runs ML-DSA-44's signing time had an interquartile range of \NarchIqrRelDSAfortyfour{} times its median and a mean \NarchTailOfDSAfortyfour{} times the median, the widest of the three, so the median timed for each held-out image, over its own set of signatures, lands further from the fixed-cost median the model uses. The next subsection shows the same tail at batch scale.

\subsection{Where the model breaks, and why}



Table~\ref{tab:archivebatch} runs 4,800 assets across worker processes and
asks whether the per-asset model predicts the batch. For ES256 it does, to about \NarchErrOneES\% at one worker. For ML-DSA-65 it is wrong by
\NarchErrOneDSAsixtyfive\% \emph{at one worker}, where there is no contention
to blame. The cause is the statistic. ML-DSA signing uses rejection sampling,
so its cost has a long right tail: at ML-DSA-65 the median fixed cost is
\NarchFixMed\,ms, the mean \NarchFixMean{} and the 95th percentile
\NarchFixPnf, a mean-to-median ratio of \NarchTail{} against \NarchTailES{} for
ES256. One signature costs about the median; a batch of 4,800 costs about the
mean. Substituting the mean cuts the error at one worker to
\NarchErrMeanOneDSAsixtyfive\%, and ES256 barely moves. Riou et al.\ warn that
rejection sampling makes ML-DSA timings heavy-tailed~\cite{riou2026}; the
batch shows the consequence for capacity. Medians suit a latency claim, and a
capacity estimate built on them for a lattice scheme is optimistic by roughly
the gap between median and mean.

Both models fail together at \NarchMaxWorkers{} workers on the M1 Pro, by \NarchErrMaxDSAsixtyfive\% and \NarchErrMeanMaxDSAsixtyfive\% for ML-DSA-65, because its six performance cores
are joined by two efficiency cores. The Pi~5's four cores are identical, so we repeated the batch there, \NarchPiRuns{} times. For ML-DSA-65 its median model errs by \NarchPiErrOneDSAsixtyfive\% at one worker (\NarchPiErrOneDSAsixtyfiveLo--\NarchPiErrOneDSAsixtyfiveHi\% across runs) and its mean model by \NarchPiErrMeanOneDSAsixtyfive\% (\NarchPiErrMeanOneDSAsixtyfiveLo--\NarchPiErrMeanOneDSAsixtyfiveHi\%), the same pattern, and at four workers the mean model errs by \NarchPiErrMeanMaxDSAsixtyfive\% (\NarchPiErrMeanMaxDSAsixtyfiveLo--\NarchPiErrMeanMaxDSAsixtyfiveHi\%) with a parallel efficiency of \NarchPiEffMaxDSAsixtyfive. On a machine without mixed cores the mean model holds for ML-DSA at every worker count, within \NarchPiErrMeanAllMaxDSAsixtyfive\% in every run. ES256 is less stable at four workers, \NarchPiErrMeanMaxESLo{} to \NarchPiErrMeanMaxESHi\% across runs, because its whole batch there takes a quarter of a second, so anything outside the per-asset model, such as starting the workers, weighs heavily.

Carried past what was measured, at the parallel efficiency observed on small
images, the model gives computation-only estimates (Table~\ref{tab:archiveproj}). For a million camera photographs re-signing adds about two
minutes to five on one machine, and for radiology series it disappears into
three hours of hashing. These leave out reading each asset from storage once,
memory bandwidth and storage traffic, which will set the real schedule for
large objects.
